\subsection{Spectrum of a compact operator}

Lemma 1. --- Let E be a separated topological vector space of dimension \(\geqslant 2\) over R, and let X be a countable subset of E. The complement E - X is connected.

First, suppose that E has dimension 2. We may assume that \(E = R^{2}\) equipped with the Euclidean norm (EVT, I, p. 14, th. 2). Since X is countable, there exists a real number \(r \in R_{+}^{*}\) such that the circle C with center 0 and radius r does not meet X. Let \(x \in E - X\) ; if \(x \notin C\) , there exists a point \(y \in C\) such that the line \(L_{x}\) joining x and y does not meet X, since X is countable. The set E−X is the union of C, which is connected, and of the connected sets \(L_{x} \cup C\) for \(x \in \mathrm{E} - (\mathrm{X} \cup \mathrm{C})\) ; these sets all contain C, and therefore E-X is connected (TG, I, p. 81, prop. 2).

Consider the general case. By replacing X with X - x for an element \(x \in E - X\) we reduce to the case where \(0 \notin X\) . Since E-X is the union of the sets \(F - (X \cap F)\) as F ranges over the set of 2-dimensional subspaces of E, and since these are connected by the preceding case and contain 0, the set E-X is connected (loc. cit.).

Lemma 2. --- Let S \(\subset \mathbf{C}\) be an infinite, discrete, bounded, and closed set in \(\mathbf{C} - \{0\}\) . Then S is the set of values of a sequence ( \(\lambda_{n}\) ) \(_{n\in N}\) of nonzero complex numbers, pairwise distinct, such that the sequence ( \(|\lambda_{n}|\) ) \(_{n\in N}\) is decreasing and converges to 0.

For every integer \(i \geqslant 1\) , the set \(\mathbf{A}_i\) of complex numbers \(\lambda \in \mathrm{S}\) such that \(|\lambda| \geqslant \frac{1}{i}\) is compact and discrete in \(\mathbf{C}\) , hence finite. Denote its cardinal by \(a_i\) . Since S is infinite, the sequence \((a_i)\) tends to \(+\infty\) . Set \(\mathbf{A}_0 = \varnothing\) and \(a_0 = 0\) . For every \(i \geqslant 1\) , let us choose a bijection \(n \mapsto \lambda_n\) from the interval \([a_{i-1}, a_i[\) of \(\mathbf{N}\) onto \(\mathbf{A}_i - \mathbf{A}_{i-1}\) such that the map \(n \mapsto |\lambda_n|\) is decreasing on \([a_{i-1}, a_i[\) . The sequence \((\lambda_n)_{n \in \mathbf{N}}\) has the required properties.

Let E be a complete normable space of infinite dimension. The algebra \(\mathcal{L}^{\mathrm{c}}(\mathrm{E})\) is a non-unital subalgebra of \(\mathcal{L}(\mathrm{E})\) . Recall that for every compact endomorphism \(u \in \mathcal{L}^{\mathrm{c}}(\mathrm{E})\) the spectrum \(\mathrm{Sp}'_{\mathcal{L}^{\mathrm{c}}(\mathrm{E})}(u)\) is the spectrum of \(u\) relative to the unital subalgebra \(\mathcal{L}^{\mathrm{c}}(\mathrm{E}) \oplus \mathbf{C1}_{\mathrm{E}}\) of \(\mathcal{L}(\mathrm{E})\) (I, p. 4, no. 4).

PROPOSITION 5. — Let E be a complete normable space and let u be a compact endomorphism of E. Every element of \(\mathrm{Sp}_{\mathrm{s}}(u)\) is an eigenvalue of finite spectral multiplicity. Moreover:

\begin{enumerate}
\def\labelenumi{\alph{enumi})}
\item
  If E is finite-dimensional, then \(\operatorname{Sp}(u) = \operatorname{Sp}_{\mathrm{s}}(u)\) ;
\item
  If E is of infinite dimension, then \(\mathrm{Sp}_{\mathrm{s}}(u) = \mathrm{Sp}(u) - \{0\}\) and \(\mathrm{Sp}_{\omega}(u) = \{0\}\) ;
\item
  If \(\mathrm{Sp}_{\mathrm{s}}(u)\) is infinite, it is the set of values of a sequence \((\lambda_{n})_{n\in\mathbf{N}}\) of nonzero complex numbers, pairwise distinct, such that the sequence \((|\lambda_{n}|)_{n\in\mathbf{N}}\) is decreasing and converges to 0;
\item
  If E is of infinite dimension, then \(\mathrm{Sp}(u) = \mathrm{Sp}'_{\mathcal{L}^c (\mathrm{E})}(u)\) .
\end{enumerate}

Every element of \(\mathrm{Sp}_{\mathrm{s}}(u)\) is an eigenvalue of finite spectral multiplicity (prop. 1 of III, p. 83).

The assertion \(a\) ) is elementary (III, p. 85, example 1). Suppose now that E is of infinite dimension.

Let \(\lambda \in \mathrm{Sp}(u) - \{0\}\) . Then \(u - \lambda 1_{\mathrm{E}}\) is a Riesz endomorphism of E (cor. 2 of prop. 2 of III, p. 75), hence \(\lambda\) belongs to \(\mathrm{Sp}_{\mathrm{s}}(u)\) . If E is infinite-dimensional, then \(\mathrm{Sp}_{\omega}(u)\) is nonempty (III, p. 87, th. 1, a)). Necessarily, \(\mathrm{Sp}_{\omega}(u) = \{0\}\) , whence \(b\) ).

The set \(\mathrm{Sp}_{\mathrm{s}}(u)\) is discrete and bounded. Moreover, we have \(\mathrm{Sp}_{\mathrm{s}}(u) = \mathrm{Sp}(u) \cap (\mathbf{C} - \{0\})\) according to \(b\) ), hence \(\mathrm{Sp}_{\mathrm{s}}(u)\) is closed in \(\mathbf{C} - \{0\}\) . The assertion \(c\) ) therefore follows from lemma 2.

The spectrum of \(u\) is countable, and its complement in \(\mathbf{C}\) is therefore connected (lemma 1). By the corollary of prop. 6 of I, p. 28, applied to the unital subalgebra \(\mathcal{L}^{\mathrm{c}}(\mathrm{E})\oplus \mathbf{C}1_{\mathrm{E}}\) of \(\mathcal{L}(\mathrm{E})\) , we conclude that \(\mathrm{Sp}(u) = \mathrm{Sp}'_{\mathcal{L}^{\mathrm{c}}(\mathrm{E})}(u)\) .

PROPOSITION 6. --- Let E be a complete normable space over C and u a compact endomorphism of E.

\begin{enumerate}
\def\labelenumi{\alph{enumi})}
\item
  Let \(f \in \mathcal{O}(\mathrm{Sp}(u))\) such that \(f(0) = 0\) . The endomorphism \(f(u)\) is compact;
\item
  Suppose that E is a complex Hilbert space and that u is normal. Let f be a continuous function on \(\mathrm{Sp}(u)\) such that \(f(0)=0\) . The normal endomorphism \(f(u)\) is compact.
\end{enumerate}

Moreover, the converse assertions are valid if E is infinite-dimensional and the condition \(f(0) = 0\) is not necessary if E is finite-dimensional.

We may assume that E is infinite-dimensional.

Let us prove \(a\) ) and its converse. The endomorphism \(u\) is an element of the Banach algebra \(\mathcal{L}^{\mathrm{c}}(\mathrm{E})\) . Since E is of infinite dimension, we have the equality \(\mathrm{Sp}(u) = \mathrm{Sp}'_{\mathcal{L}^{\mathrm{c}}(\mathrm{E})}(u)\) by prop. 5, \(d\) ). The element \(f(u)\) of the holomorphic functional calculus of the unital Banach algebra derived from \(\mathcal{L}^{\mathrm{c}}(\mathrm{E})\) by adjoining an identity element belongs to \(\mathcal{L}^{\mathrm{c}}(\mathrm{E})\) if and only if \(f(0) = 0\) (I, p. 88). But moreover this element coincides with the element \(f(u)\) of \(\mathcal{L}(\mathrm{E})\) (prop. 7 of I, p. 75), hence \(f(u)\) is compact if and only if \(f(0) = 0\) .

The proof of assertion \(b\) ) and of its converse is entirely identical, considering the continuous functional calculus of the C*-algebra \(\mathcal{L}^{\mathrm{c}}(\mathrm{E})\) (I, p. 110, def. 5).

COROLLARY. --- Let E and F be Hilbert spaces and let u be a continuous linear map from E into F. The linear map u is compact if and only if the endomorphism \(|u|\) of E is compact.

Let \((j, |u|)\) be the polar decomposition of u (def. 4 of I, p. 140). Since \(u = j|u|\) and \(|u| = \sqrt{u^{*}u}\) (prop. 10 of I, p. 139), the equivalence follows from prop. 3 of III, p. 5 and from assertion b) of the preceding proposition.

